% GENERATED FILE. Edit canonical lessons, then run npm run generate:books.
% canonical-source-hash: fnv1a64:3e39df33c6c0fcba
% canonical-lessons: PT-C415-carruagem, PT-C415-compartimento, PT-C415-beliche, PT-C415-maquinista, PT-C415-locomotiva

\chapter{No comboio}
\label{ch:pt-no-comboio}
\hlchaptermodality{415}

\begin{chapteropening}
\textbf{By the end of this chapter:} \emph{I can find my place on a train: the carriage, the compartment, the bunk, the train driver and the engine.}

Complete the last lesson of chapter 415: I can find my place on a train: the carriage, the compartment, the bunk, the train driver and the engine.
\end{chapteropening}

\section[a carruagem]{a carruagem — the carriage (of a train); the coach (horse-drawn)}

\label{lesson:PT-C415-carruagem}



\begin{quote}
Before the new one: say the Portuguese for the ticket inspector, then the Portuguese for the fine.
\end{quote}

\subsection*{You'll want to know: a carruagem}
\textbf{a carruagem} — \textquotedblleft{}the carriage (of a train); the coach (horse-drawn)\textquotedblright{}.

\begin{quote}
O nosso lugar é na carruagem cinco.
\end{quote}

\begin{quote}
\emph{Our seat is in carriage five.}
\end{quote}

In Brazil a train carriage is \emph{o vagão}.

\begin{cousinweb}
\emph{carruagem} came from Spanish \emph{carruaje}, which came from Catalan \emph{carruatge} (EtymDB); the family goes back to Latin \emph{carrus} 'cart'.
\end{cousinweb}

\begin{grammarlens}[title={What you've built}]
One more word for telling, explaining and describing at B1.
\end{grammarlens}

\subsection*{Your turn}
\begin{itemize}
\raggedright
  \item \emph{Say it:} \emph{a carruagem}
  \item \emph{Say it:} \emph{a carruagem}, once more
  \item \emph{Say it:} a sentence of your own with \emph{a carruagem}
  \item \emph{Recall:} say the Portuguese for the ticket inspector, then the Portuguese for the fine, then say \emph{a carruagem} again
\end{itemize}

\begin{tcolorbox}[breakable,colback=teal!4,colframe=teal!35!black,title={Before you move on}]
What does a carruagem mean? (\textquotedblleft{}The carriage (of a train); the coach (horse-drawn)\textquotedblright{}.) Say it once more.
\end{tcolorbox}
\section[o compartimento]{o compartimento — the compartment}

\label{lesson:PT-C415-compartimento}



\begin{quote}
Before the new one: say the Portuguese for the fine, then the Portuguese for the carriage.
\end{quote}

\subsection*{You'll want to know: o compartimento}
\textbf{o compartimento} — \textquotedblleft{}the compartment\textquotedblright{}.

\begin{quote}
Partilhamos o compartimento com uma família simpática.
\end{quote}

\begin{quote}
\emph{We share the compartment with a nice family.}
\end{quote}

Also a compartment in a bag or a drawer.

\begin{cousinweb}
\emph{compartimento} is a learned formation on Late Latin \emph{compartīrī} 'to share out', built on \emph{pars} 'part' (Lewis \& Short).
\end{cousinweb}

\begin{grammarlens}[title={What you've built}]
One more word for telling, explaining and describing at B1.
\end{grammarlens}

\subsection*{Your turn}
\begin{itemize}
\raggedright
  \item \emph{Say it:} \emph{o compartimento}
  \item \emph{Say it:} \emph{o compartimento}, once more
  \item \emph{Say it:} a sentence of your own with \emph{o compartimento}
  \item \emph{Recall:} say the Portuguese for the fine, then the Portuguese for the carriage, then say \emph{o compartimento} again
\end{itemize}

\begin{tcolorbox}[breakable,colback=teal!4,colframe=teal!35!black,title={Before you move on}]
What does o compartimento mean? (\textquotedblleft{}The compartment\textquotedblright{}.) Say it once more.
\end{tcolorbox}
\section[o beliche]{o beliche — the bunk (bed)}

\label{lesson:PT-C415-beliche}



\begin{quote}
Before the new one: say the Portuguese for the carriage, then the Portuguese for the compartment.
\end{quote}

\subsection*{You'll want to know: o beliche}
\textbf{o beliche} — \textquotedblleft{}the bunk (bed)\textquotedblright{}.

\begin{quote}
No comboio da noite, durmo no beliche de cima.
\end{quote}

\begin{quote}
\emph{On the night train, I sleep in the top bunk.}
\end{quote}

At home, two children often share \emph{um beliche}.

\begin{cousinweb}
The origin of \emph{beliche} is not settled. It first meant a sleeping place on a ship, and EtymDB has no entry for it.
\end{cousinweb}

\begin{grammarlens}[title={What you've built}]
One more word for telling, explaining and describing at B1.
\end{grammarlens}

\subsection*{Your turn}
\begin{itemize}
\raggedright
  \item \emph{Say it:} \emph{o beliche}
  \item \emph{Say it:} \emph{o beliche}, once more
  \item \emph{Say it:} a sentence of your own with \emph{o beliche}
  \item \emph{Recall:} say the Portuguese for the carriage, then the Portuguese for the compartment, then say \emph{o beliche} again
\end{itemize}

\begin{tcolorbox}[breakable,colback=teal!4,colframe=teal!35!black,title={Before you move on}]
What does o beliche mean? (\textquotedblleft{}The bunk (bed)\textquotedblright{}.) Say it once more.
\end{tcolorbox}
\section[o maquinista]{o maquinista — the train driver}

\label{lesson:PT-C415-maquinista}



\begin{quote}
Before the new one: say the Portuguese for the compartment, then the Portuguese for the bunk.
\end{quote}

\subsection*{You'll want to know: o maquinista}
\textbf{o maquinista} — \textquotedblleft{}the train driver\textquotedblright{}.

\begin{quote}
O maquinista avisou que o comboio ia parar.
\end{quote}

\begin{quote}
\emph{The train driver announced that the train was going to stop.}
\end{quote}

A bus driver is \emph{o motorista}.

\begin{cousinweb}
\emph{maquinista} is built on \emph{máquina} 'machine, engine' with \emph{-ista}; \emph{máquina} is a learned borrowing from Latin \emph{māchina} 'machine, device' (Lewis \& Short).
\end{cousinweb}

\begin{grammarlens}[title={What you've built}]
One more word for telling, explaining and describing at B1.
\end{grammarlens}

\subsection*{Your turn}
\begin{itemize}
\raggedright
  \item \emph{Say it:} \emph{o maquinista}
  \item \emph{Say it:} \emph{o maquinista}, once more
  \item \emph{Say it:} a sentence of your own with \emph{o maquinista}
  \item \emph{Recall:} say the Portuguese for the compartment, then the Portuguese for the bunk, then say \emph{o maquinista} again
\end{itemize}

\begin{tcolorbox}[breakable,colback=teal!4,colframe=teal!35!black,title={Before you move on}]
What does o maquinista mean? (\textquotedblleft{}The train driver\textquotedblright{}.) Say it once more.
\end{tcolorbox}
\section[a locomotiva]{a locomotiva — the engine, the locomotive}

\label{lesson:PT-C415-locomotiva}



\begin{quote}
Before the new one: say the Portuguese for the bunk, then the Portuguese for the train driver.
\end{quote}

\subsection*{You'll want to know: a locomotiva}
\textbf{a locomotiva} — \textquotedblleft{}the engine, the locomotive\textquotedblright{}.

\begin{quote}
A locomotiva antiga está agora no museu.
\end{quote}

\begin{quote}
\emph{The old engine is now in the museum.}
\end{quote}

An older word: \emph{a máquina} of the train.

\begin{cousinweb}
\emph{locomotiva} is a learned coinage from Latin \emph{locō} 'from a place' and Medieval Latin \emph{mōtīvus} 'moving', from \emph{movēre} 'to move' (Lewis \& Short has \emph{movēre}).
\end{cousinweb}

\begin{grammarlens}[title={What you've built}]
One more word for telling, explaining and describing at B1.
\end{grammarlens}

\subsection*{Your turn}
\begin{itemize}
\raggedright
  \item \emph{Say it:} \emph{a locomotiva}
  \item \emph{Say it:} \emph{a locomotiva}, once more
  \item \emph{Say it:} a sentence of your own with \emph{a locomotiva}
  \item \emph{Recall:} say the Portuguese for the bunk, then the Portuguese for the train driver, then say \emph{a locomotiva} again
\end{itemize}

\begin{tcolorbox}[breakable,colback=teal!4,colframe=teal!35!black,title={Before you move on}]
What does a locomotiva mean? (\textquotedblleft{}The engine, the locomotive\textquotedblright{}.) Say it once more.
\end{tcolorbox}
